\documentclass[11pt]{article}
\usepackage{lmodern}
\usepackage[T1]{fontenc}
\usepackage{amsmath,amssymb,amsthm}
\usepackage[margin=1in]{geometry}
\usepackage[hidelinks]{hyperref}
\title{Fourier Residues of a Plane Cubic, Rational Chords, and the Elliptic L-Function}
\author{\small\href{mailto:admin@cosmosgeek.org}{admin@cosmosgeek.org}}
\date{Working draft for specialist review}
\begin{document}
\maketitle

\begin{abstract}
We study a prime-power Fourier sum for the affine cone over an elliptic
curve embedded as a plane cubic. Its double-pole coefficient at a
primitive frequency is nonzero precisely when the corresponding line
cuts out three rational points, counted with multiplicity. The weights
for ordinary chords, tangents and flexes are in the ratio $6:3:1$,
including at primes where distinct rational points have the same
reduction. The resulting measure is a constant multiple of the
pushforward of counting measure on pairs of rational points. Its
height exponent is the Mordell--Weil rank. On the multiples of any
one split frequency, a Dirichlet series of the same coefficients
recovers the analytic rank by meromorphic continuation. We also
prove uniformity for growing transverse chord height and a finite
measure criterion for arithmetic rank zero. On a pencil through a nonflex
rational point we compute an all-prime boundary coefficient whose continuation
has a half-integral singular exponent in odd rank. The equality of the two
rank exponents remains open.
\end{abstract}

\section{The cubic and its frequencies}

Let $E/\mathbb Q$ be an elliptic curve embedded in
$\mathbb P^2$ by the complete linear system $|3O|$,
with primitive integral plane-cubic equation, smooth over $\mathbb Q$,
\[
F(X_1,X_2,X_3)=0.
\]
Write
\[
r=\operatorname{rank}E(\mathbb Q),\qquad
a=\operatorname{ord}_{s=1}L(E,s).
\]
Let $S$ be a finite nonempty set containing
$2,3,5$ and every prime of bad reduction of
the chosen integral model.

All Euler products below are over $p\notin S$.
Chambert-Loir and Tschinkel studied the blowup of the ideal $(x_0,F)$
and its local additive Fourier transforms \cite[\S3]{CLT}. We work
with explicitly normalized finite exponential sums at the retained primes.
The archimedean test function and the local factors at $S$ are omitted.
Write $\zeta^S(s)=\zeta(s)\prod_{p\in S}(1-p^{-s})$ and use $L^S$
for the elliptic $L$-function with its factors at $S$ removed.

For a nonzero integral frequency
$\xi=(\xi_1,\xi_2,\xi_3)$, write
\[
L_\xi:\quad \xi_1X_1+\xi_2X_2+\xi_3X_3=0.
\]
Call $\xi$ primitive if
$\gcd(\xi_1,\xi_2,\xi_3)=1$.
Primitive vectors differing by sign represent
the same rational projective line.

\subsection{The exact local shell factor}

For $p\notin S$, put
\[
n_p=\#E(\mathbb F_p)=p+1-a_p.
\]
For $k\geq1$ define
\[
R_{p,k}(\xi)=
\#\left\{
P\in E(\mathbb Z/p^k):
\xi(P)\equiv0\pmod{p^k}
\right\}.
\]
The condition is independent of the choice of
primitive homogeneous coordinates for $P$.

Write $e_m(t)=\exp(2\pi it/m)$. Our Fourier coefficient is defined by
\[
 S_{p^k}(\xi)=\sum_{\substack{x\in(\mathbb Z/p^k)^3,\ x\not\equiv0\pmod p\\
 F(x)\equiv0\pmod{p^k}}} e_{p^k}(\xi\cdot x).
\]
For example, at a general modulus $d$ prime to $S$ these coefficients
are the Fourier coefficients of the counting measure at the fractions
$x/d$ on $\mathbb R^3/\mathbb Z^3$, with the same primitivity and cubic
congruence imposed modulo $d$. The Chinese remainder theorem factors
the corresponding sum into the prime-power sums above. The inverses
introduced in the additive characters are units, which permute the
allowed scalar multiples on each cone.

Summing first over the unit scalar multiplying
a projective point gives the Ramanujan sums
\[
c_{p^k}(n)=
\begin{cases}
p^k-p^{k-1},&p^k\mid n,\\
-p^{k-1},&p^{k-1}\mid n,\ p^k\nmid n,\\
0,&p^{k-1}\nmid n.
\end{cases}
\]
Since the smooth curve has $p$ lifts of each
point from modulus $p^{k-1}$ to modulus $p^k$,
we obtain
\begin{align}
S_p(\xi)&=pR_{p,1}(\xi)-n_p,
\label{eq:ram1}\\
S_{p^k}(\xi)
&=p^k\bigl(R_{p,k}(\xi)-R_{p,k-1}(\xi)\bigr),
\quad k\geq2.
\label{eq:ramk}
\end{align}

For $\Re w>2$ define
\[
F_{p,\xi}(w)=
1+\sum_{k\geq1}p^{-kw}S_{p^k}(\xi).
\]
Writing $w=2+u$, equations
\eqref{eq:ram1}--\eqref{eq:ramk} give
\begin{equation}\label{eq:local-R}
F_{p,\xi}(2+u)
=
1-n_pp^{-2-u}
+(1-p^{-1-u})
\sum_{k\geq1}R_{p,k}(\xi)p^{-k(1+u)}.
\end{equation}

This normalization counts every primitive scalar representative.
The use of hyperplane-section counts at prime powers is also present
in the local integrals of \cite[\S3, Lemma~3.3]{CLT}.

For the zero frequency, smooth Hensel lifting
gives
\[
S_{p^k}(0)=(p-1)n_pp^{2k-2}.
\]
Thus
\begin{equation}\label{eq:zero-local}
F_{p,0}(2+u)
=
1+c_p\sum_{k\geq1}p^{-ku},
\qquad
c_p=\frac{(p-1)n_p}{p^2}.
\end{equation}

\subsection{A constant local value, including collisions}

Suppose that $L_\xi$ has a completely split
intersection divisor over $\mathbb Q$:
\begin{equation}\label{eq:split-divisor}
(L_\xi\cdot E)_{\mathbb Q}
=[P_1]+[P_2]+[P_3],
\qquad P_i\in E(\mathbb Q),
\end{equation}
with repetitions permitted. In this and the next two subsections the frequency is primitive.

At every good prime $p$, the points $P_i$
extend to sections of the smooth proper
model $\mathcal E/\mathbb Z_p$.

Because $\xi$ is primitive, the reduction
of $L_\xi$ does not contain the entire
special fibre. Therefore its horizontal
Cartier divisor on $\mathcal E$ is exactly
the sum of the three sections in
\eqref{eq:split-divisor}, with
multiplicities and no vertical component.

Let $\mu_p$ denote the translation-invariant
measure on $E(\mathbb Q_p)$ normalized so
that every residue class modulo $p^k$
has measure $p^{-k}$.
Its total mass is $n_p/p$.

For the line section put
\[
V_\xi(P)=v_p(\xi(P)).
\]
The divisor identity gives
\[
V_\xi(P)=V_{P_1}(P)+V_{P_2}(P)+V_{P_3}(P),
\]
where $V_{P_i}$ is the valuation of a
local parameter cutting out the section
$P_i$.

For each section,
\[
\int V_{P_i}\,d\mu_p
=\sum_{k\geq1}p^{-k}
=\frac1{p-1}.
\]
This remains true when two or three sections
coincide modulo arbitrarily high powers of $p$.
Additivity therefore gives
\[
\sum_{k\geq1}R_{p,k}(\xi)p^{-k}
=\int V_\xi\,d\mu_p
=\frac3{p-1}.
\]

Substituting into \eqref{eq:local-R} at $u=0$
proves the exact identity
\begin{equation}\label{eq:constant-local}
\boxed{
F_{p,\xi}(2)
=
1-\frac{n_p}{p^2}+\frac3p
=:F_p^*(0).
}
\end{equation}

In particular, the value is independent
of the rational split line and of all
collisions between its three sections.

Since $n_p=p+1-a_p$,
\[
F_p^*(0)
=1+\frac2p+\frac{a_p-1}{p^2}>0
\]
for the retained primes.

\subsection{Pole order and orbifold multiplicities}

Define the good-prime Euler product
\[
D_\xi^S(2+u)
=\prod_{p\notin S}F_{p,\xi}(2+u).
\]

For a transverse rational split line,
away from finitely many collision primes,
\[
R_{p,k}(\xi)=3\quad(k\geq1).
\]
Hence
\[
F_{p,\xi}(2+u)
=
1-n_pp^{-2-u}+3p^{-1-u}
=
1+2p^{-1-u}+O(p^{-3/2})
\]
uniformly in a sufficiently small
right-hand neighborhood of $u=0$.

Thus the generic factor is
$\zeta^S(1+u)^2$.

For a rational tangent line with divisor
$2[P]+[Q]$, $P\ne Q$, the generic local
intersection is a unit times
$t^2$ at $P$ and a unit times $t$ at $Q$.
Consequently
\[
R_{p,k}=1+p^{\lfloor k/2\rfloor}
\]
outside finitely many exceptional primes.
The first two nonuniform shell contributions
give
\[
F_{p,\xi}(2+u)
=
1+p^{-1-u}+p^{-1-2u}
+O(p^{-3/2+\epsilon}).
\]
Its two zeta factors are therefore
\[
\zeta^S(1+u)\zeta^S(1+2u).
\]

For a rational flex line with divisor
$3[P]$, the generic local equation is
a unit times $t^3$. Then
\[
R_{p,k}=p^{k-\lceil k/3\rceil},
\]
and
\[
F_{p,\xi}(2+u)
=
1+p^{-1-2u}+p^{-1-3u}
+O(p^{-3/2+\epsilon}).
\]
Its two factors are
\[
\zeta^S(1+2u)\zeta^S(1+3u).
\]

All error terms above become absolutely
summable Euler-product corrections in a
small neighborhood of $u=0$.

Put
\[
\rho_S=\prod_{p\in S}(1-p^{-1}),
\]
and
\[
A_E^S=
\prod_{p\notin S}
(1-p^{-1})^2F_p^*(0).
\]
By Hasse's bound,
\[
(1-p^{-1})^2F_p^*(0)
=1+O(p^{-3/2}),
\]
so $A_E^S$ converges absolutely
to a positive number.

Define
\begin{equation}\label{eq:CE}
C_E^S=\rho_S^2A_E^S>0.
\end{equation}
The factor $\rho_S^2$ is essential:
$\zeta^S(1+u)\sim\rho_S/u$,
not $1/u$.

For each rational split projective line,
the limit
\[
R_\xi^S=
\lim_{u\downarrow0}u^2D_\xi^S(2+u)
\]
therefore exists and equals
\begin{equation}\label{eq:orbifold}
\boxed{
R_\xi^S=
\begin{cases}
C_E^S,&P_1,P_2,P_3\text{ distinct},\\[2mm]
C_E^S/2,&L_\xi\cdot E=2P+Q,\ P\ne Q,\\[2mm]
C_E^S/6,&L_\xi\cdot E=3P.
\end{cases}}
\end{equation}

The constant is the same in all cases:
the regularized Euler product at $u=0$
depends on the values \eqref{eq:constant-local},
which are universal. Only the slopes
$1,2,3$ of the zeta factors change.

For a reduced intersection divisor that is not completely split over
$\mathbb Q$, let $A_\xi$ be its degree-three finite \'etale algebra.
Outside its finite exceptional set, $R_{p,k}$ is constant in $k$ and
equals the number of fixed points of Frobenius on its three geometric
points. The first coefficient agrees with that of
$\zeta_{A_\xi}(1+u)/\zeta(1+u)$; the quotient of Euler factors is
$1+O(p^{-3/2+3\epsilon})$ near $u=0$ for small $\epsilon>0$.
At the finitely many remaining good primes the reduced cubic has
distinct geometric roots in characteristic zero. On disjoint
$p$-adic neighborhoods of its rational roots its equation is a local
parameter times a factor of fixed valuation, and away from those
neighborhoods the valuation is bounded by compactness. Thus $R_{p,k}$
is bounded for large $k$. The factors in \eqref{eq:local-R} are regular
and positive at $u=0$ since $n_p<p^2$. Each field factor of $A_\xi$
has a Dedekind zeta function with a simple pole at one. The resulting
pole order is
\[
\#\{\text{field factors of }A_\xi\}-1<2.
\]
Hence its $u^2$ residue is zero.
The nonreduced intersections are the
rational tangent and flex cases above.

\subsection{Exact chord-map measure}

Let
\[
\Phi:E^2\longrightarrow(\mathbb P^2)^\vee
\]
be the chord map. It sends $(P,Q)$
to the line through $P,Q$ when $P\ne Q$,
and to the tangent at $P$ when $P=Q$.
The third intersection is
\[
R=-P-Q.
\]

The action of $S_3$ on ordered triples
$(P,Q,R)$ with $P+Q+R=O$
makes $\Phi$ the finite degree-six quotient
\[
E^2/S_3\simeq |3O|\simeq\mathbb P^2.
\]

A transverse split line has six ordered
preimages; a tangent line $2P+Q$
has three; a flex $3P$ has one.
Thus \eqref{eq:orbifold} gives the
following exact identity of atomic
measures on rational projective lines:
\begin{equation}\label{eq:measure-push}
\boxed{
\sum_{\xi\ {\rm rational\ split}}
R_\xi^S\,\delta_\xi
=
\frac{C_E^S}{6}\,
\Phi_*\left(
\sum_{(P,Q)\in E(\mathbb Q)^2}
\delta_{(P,Q)}
\right).
}
\end{equation}

This includes the flex at $O$ and applies
when the Mordell--Weil rank is zero.
It is an identity for the finite Euler
residues, not for the complete adelic
Fourier transforms.

\subsection{The height exponent}

Let $H(\xi)$ be the standard primitive
projective height on the dual plane.
Choose a basis of $E(\mathbb Q)$ modulo torsion. Write $G$ for the
Gram matrix of the N\'eron--Tate height attached to $\mathcal O_E(O)$,
normalized by $\hat h(\sum_i n_i P_i)=n^{\mathrm T}Gn$. Write
\[
\operatorname{Reg}(E)=\det G,
\qquad T=E(\mathbb Q)_{\rm tors}.
\]

To determine the chord-map height,
choose the hyperplane in $|3O|$
consisting of divisors containing $O$.
Its pullback under $\Phi$ is
\[
\{P=O\}+\{Q=O\}+\{P+Q=O\}.
\]
Consequently
\[
\Phi^*\mathcal O(1)\simeq
p_1^*\mathcal O_E(O)
\otimes p_2^*\mathcal O_E(O)
\otimes m^*\mathcal O_E(O),
\]
where $m(P,Q)=P+Q$.

The associated canonical quadratic height is
\[
Q(P,Q)=
\hat h(P)+\hat h(Q)+\hat h(P+Q).
\]
The ordinary projective height differs by
a bounded term:
\begin{equation}\label{eq:chord-height}
\log H(\Phi(P,Q))=Q(P,Q)+O_E(1).
\end{equation}

On the free part of $E(\mathbb Q)^2$,
the height matrix is
\[
\begin{pmatrix}2&1\\1&2\end{pmatrix}
\otimes G,
\]
whose determinant is
\[
3^r\operatorname{Reg}(E)^2.
\]

Geometry of numbers in dimension $2r$
therefore gives, for $r>0$,
\begin{equation}\label{eq:pair-count}
\#\{(P,Q)\in E(\mathbb Q)^2:
H(\Phi(P,Q))\le B\}
\sim
\frac{|T|^2\pi^r}
{r!\,3^{r/2}\operatorname{Reg}(E)}
(\log B)^r.
\end{equation}

For $r=0$, the left side is eventually
the constant $|T|^2$; the same formula
holds with $\operatorname{Reg}(E)=1$
and $(\log B)^0=1$.

Combining \eqref{eq:measure-push}
and \eqref{eq:pair-count} proves
\begin{equation}\label{eq:weighted-count}
\boxed{
\sum_{\substack{\xi\ {\rm rational\ split}\\
H(\xi)\le B}}R_\xi^S
\sim
\frac{C_E^S|T|^2\pi^r}
{6r!\,3^{r/2}\operatorname{Reg}(E)}
(\log B)^r.
}
\end{equation}

Thus the primitive split residue measure
recovers the Mordell--Weil rank through
the exponent of logarithmic height growth.

\subsection{Uniformity and a joint split-sector limit}

For a transverse split line, let
$\Delta_\xi\ne0$ denote its integral
dual-sextic discriminant.
Outside the primes dividing $\Delta_\xi$,
the three intersection sections remain
distinct and
\[
F_{p,\xi}(2+u)
=
F_p^*(u):=
1-n_pp^{-2-u}+3p^{-1-u}.
\]

At a collision prime, write
\[
V_\xi=V_{P_1}+V_{P_2}+V_{P_3}
\]
as above.
The exact second moment of an individual
section valuation is
\[
\int V_{P_i}^2\,d\mu_p
=
\sum_{k\ge1}(2k-1)p^{-k}
=
\frac{p+1}{(p-1)^2}.
\]
By Cauchy--Schwarz,
\[
\int V_\xi^2\,d\mu_p
\le3\sum_i\int V_{P_i}^2\,d\mu_p
\ll p^{-1},
\]
uniformly in the collision depth.

Equation \eqref{eq:local-R}, expressed as
\[
F_{p,\xi}(2+u)
=
1-n_pp^{-2-u}
+(1-p^{-1-u})
\int\sum_{k=1}^{V_\xi(P)}
p^{-uk}\,d\mu_p(P),
\]
therefore gives
\[
\left|
\frac{\partial}{\partial u}
\bigl(F_{p,\xi}(2+u)-F_p^*(u)\bigr)
\right|
\ll\frac{\log p}{p},
\qquad 0\le u\le u_0.
\]

The difference vanishes at $u=0$
by \eqref{eq:constant-local}.
For the retained primes the factors
are uniformly bounded away from zero.
Hence
\[
\log
\frac{D_\xi^S(2+u)}
{\prod_{p\notin S}F_p^*(u)}
\ll
u\sum_{p\mid\Delta_\xi}\frac{\log p}{p}.
\]

The elementary maximal-order estimate
\[
\sum_{p\mid n}\frac{\log p}{p}
\ll\log\log(3|n|)
\]
and the degree-six bound
\[
|\Delta_\xi|\ll_E H(\xi)^6
\]
yield
\begin{equation}\label{eq:uniform-ratio}
\boxed{
\frac{D_\xi^S(2+u)}
{\prod_{p\notin S}F_p^*(u)}
=
1+O_E\bigl(u\log(1/u)\bigr)
}
\end{equation}
uniformly for transverse split lines with
\[
H(\xi)\le e^{1/u}.
\]

The universal product satisfies
\[
u^2\prod_{p\notin S}F_p^*(u)
\longrightarrow C_E^S.
\]
Therefore, for $r>0$, the transverse
split sector has the genuine joint limit
\begin{equation}\label{eq:joint-limit}
\boxed{
\begin{aligned}
\lim_{u\downarrow0}
u^{r+2}
\sum_{\substack{\xi\ {\rm primitive,\ split,}\\
\text{transverse}\\H(\xi)\le e^{1/u}}}
D_\xi^S(2+u)
=
\frac{C_E^S|T|^2\pi^r}
{6r!\,3^{r/2}\operatorname{Reg}(E)}.
\end{aligned}
}
\end{equation}

In deriving the transverse count from \eqref{eq:pair-count},
repeated points contribute $O((\log B)^{r/2})$ ordered pairs:
on a diagonal the quadratic height is $6\hat h(P)$. The other two
diagonals are its permutations. The flexes belong to $E[3]$.

When $r=0$ the limit for the transverse lines alone is $C_E^S$
times their finite number. Including all split lines gives the total
$C_E^S|T|^2/6$ by \eqref{eq:measure-push}.

The uniform estimate is essential:
\eqref{eq:joint-limit} does not follow
from exchanging individual residues
with an infinite frequency sum.

\subsection{Content and the original elliptic L-function}

Fix any primitive rational split frequency
$\xi$, allowing transverse, tangent or flex
contact. Let $h\ge1$ be prime to $S$.

If $p^t\parallel h$, then the first $t$
additive conductor shells of $h\xi$
have zero frequency.
For later shells, reduction by $p^t$
and smooth lifting give
\[
S_{p^k}(h\xi)=p^{2t}S_{p^{k-t}}(\xi)
\qquad(k>t).
\]
Using \eqref{eq:zero-local},
\[
\boxed{
F_{p,h\xi}(2+u)
=
(1-p^{-tu})F_{p,0}(2+u)
+p^{-tu}F_{p,\xi}(2+u).
}
\]

At $u=0$ this yields
\[
F_{p,h\xi}(2)
=
F_p^*(0)+t c_p.
\]
Thus the residue changes only by the
content:
\begin{equation}\label{eq:content-residue}
\boxed{
R_{h\xi}^S
=
R_\xi^S
\prod_{p^t\parallel h}
\left(1+\frac{t c_p}{F_p^*(0)}\right).
}
\end{equation}

In particular, the content multiplier is
independent of the split direction.

For $\Re z>1$, define
\[
\mathcal C_\xi^S(z)=
\sum_{\substack{h\ge1\\(h,S)=1}}
\frac{R_{h\xi}^S}{h^z}.
\]
Put
\[
\lambda_p=
\frac{a_p+2}{pF_p^*(0)}.
\]
The exact identity
\[
F_p^*(0)-c_p=\frac{a_p+2}{p}
\]
gives
\begin{align}
\mathcal C_\xi^S(z)
&=
R_\xi^S
\prod_{p\notin S}
\sum_{t\ge0}
\left(1+t\frac{c_p}{F_p^*(0)}\right)
p^{-tz}\notag\\
&=
\boxed{
R_\xi^S\,\zeta^S(z)^2
\prod_{p\notin S}(1-\lambda_pp^{-z}).
}
\label{eq:content-euler}
\end{align}

For the original elliptic curve, the good
Euler denominator is
\[
L_p(E,1+z)^{-1}
=
1-a_pp^{-1-z}+p^{-1-2z}.
\]
Define
\[
V_p(z)=
\frac{1-\lambda_pp^{-z}}
{(1-a_pp^{-1-z}+p^{-1-2z})
 (1-p^{-1-z})^2(1-p^{-1-2z})}.
\]

Using $|a_p|\le2\sqrt p$,
the linear and quadratic logarithmic
terms cancel, giving
\[
V_p(z)=1+O(p^{-3/2+3\epsilon})
\]
locally uniformly for $\Re z\ge-\epsilon$,
with any fixed $\epsilon<1/6$,
for all sufficiently large retained primes. For the finitely many others
the factors and their denominators are nonzero at $z=0$ as the next
formula shows.

Hence
\[
V_E^S(z)=\prod_{p\notin S}V_p(z)
\]
is holomorphic and nonzero near $z=0$.

At $z=0$,
\[
V_p(0)
=
\frac1{(1-p^{-1})^2F_p^*(0)},
\]
so the normalization is exact:
\[
\boxed{V_E^S(0)=(A_E^S)^{-1}.}
\]

We obtain the meromorphic identity
\begin{equation}\label{eq:content-L}
\boxed{
\mathcal C_\xi^S(z)=
\frac{
R_\xi^S\zeta^S(z)^2V_E^S(z)}
{
L^S(E,1+z)\,
\zeta^S(1+z)^2\,
\zeta^S(1+2z)
}.
}
\end{equation}

Write
\[
a=\operatorname{ord}_{s=1}L(E,s).
\]
The incomplete factors at $1+z$
have their usual poles.
But because $S$ is nonempty,
\[
\zeta^S(z)
=
\zeta(z)\prod_{p\in S}(1-p^{-z})
\]
has a zero of order $|S|$ at $z=0$.

Consequently the physical content series
in \eqref{eq:content-L} has order
\[
\boxed{
\operatorname{ord}_{z=0}
\mathcal C_\xi^S(z)=2|S|+3-a.
}
\]

If the known elementary omissions are
restored by defining
\[
\widetilde{\mathcal C}_\xi^S(z)
=
\mathcal C_\xi^S(z)
\prod_{p\in S}(1-p^{-z})^{-2},
\]
then
\begin{equation}\label{eq:analytic-order}
\boxed{
\operatorname{ord}_{z=0}
\widetilde{\mathcal C}_\xi^S(z)=3-a.
}
\end{equation}

The restoration is a meromorphic
normalization, not the addition of
previously counted frequencies.

Thus the analytic rank is recovered from
the order of the content Dirichlet series
attached to any single rational split
frequency, after the explicit finite-set
correction.

\section{A norm for the primitive directions}

For an integral frequency $\eta=hv\ne0$ with $v$ primitive and $h>0$,
extend the content multiplier $g(h)$ of \eqref{eq:content-residue}
by ignoring primes in $S$. Those primes do not occur in $D_\eta^S$.
After division by $g(h)$ the residues have a value $b(hv)=R_v^S$
constant along each rational line, and are zero along nonsplit lines.

A direction of primitive height $H(v)$ has $O(B/H(v))$ nonzero
multiples in the ball $\|\eta\|_\infty\le B$. By \eqref{eq:weighted-count}
the sum of $1/H(v)$ over distinct split directions converges, since
each nonzero weight is at least $C_E^S/6$. The number of supported
integral frequencies in that ball is consequently $O_E(B)$.

On the three-torus let $T_v=\{x:v\cdot x=0\bmod1\}$, and let $m_v$
be its Haar probability measure. The Fourier transform of $m_v$ is
the indicator of $\mathbb Z v$. Retain the \emph{whole} lines of
primitive height at most $B$ and set all the other nonzero coefficients
to zero. The minimum total variation of a finite measure with these
Fourier coefficients, allowing any coefficient at zero, is exactly
\[
 \sum_{\substack{[v]\ {\rm split}\\H(v)\le B}}R_v^S .
\]
Indeed, $\sum_v R_v^S m_v$ realizes these coefficients. Every other
realization differs by a multiple of Haar measure on the three-torus,
by uniqueness of Fourier coefficients. This Haar measure is singular
to the finite sum of subtorus measures. The minimum therefore has the
height asymptotic in \eqref{eq:weighted-count}.

For the full set of directions a finite complex measure with these
coefficients exists exactly when $r=0$. The finite sum just constructed
settles that case. Conversely, for a putative measure $\mu$ the averages
of its coefficients at $jv$, $|j|\le n$, give $\mu(T_v)=R_v^S$ by
dominated convergence. For two distinct split directions $v,v'$,
average instead at $\eta_0+jv+kv'$, $|j|,|k|\le n$. For fixed $\eta_0$
these are distinct and lie in a ball of radius $O(n+1)$. Only $O(n+1)$
coefficients are nonzero and all are bounded, so the limit is zero.
The same dominated convergence argument identifies this limit with
the $\eta_0$-th Fourier coefficient of $\mu$ restricted to
$T_v\cap T_{v'}$. Uniqueness gives $|\mu|(T_v\cap T_{v'})=0$.
Thus $\|\mu\|$ is at least the sum of $R_v^S$ over any finite list of
distinct split lines. For $r>0$ these sums diverge by
\eqref{eq:weighted-count}.

\section{The comparison left open}

Equations \eqref{eq:weighted-count} and
\eqref{eq:analytic-order} recover two
different integers from one family of
finite-place Fourier residues:
\[
r=\operatorname{rank}E(\mathbb Q)
\]
from the logarithmic growth of the
primitive split support, and
\[
a=\operatorname{ord}_{s=1}L(E,s)
\]
from the meromorphic content factor.

The radial scaling of an integral
frequency satisfies
\[
\|h\xi\|_\infty=h\|\xi\|_\infty,
\]
so its content Mellin parameter is
compatible with the ordinary frequency
height. This does not identify the two
rank exponents.

The norm above is defined by retaining entire primitive directions.
A global Poisson comparison would have to control the archimedean
transform and the meromorphic continuation while the frequency range
grows. A finite cutoff with zero mean removes the zero Fourier mode
altogether.

The missing implication is a rank-independent identity or bound
forcing the growth exponent in \eqref{eq:weighted-count} to equal
$3-\operatorname{ord}_{z=0}\widetilde{\mathcal C}_\xi^S(z)$.
Both quantities are determined by the same residue family, but the
formulas evaluate them by different operations. Their equality is
the rank conjecture in the official statement \cite{Clay}.

\appendix
\section{A one-sided coarea--Euler residue on a rational pencil}
\label{app:coarea-euler}

A pencil is the family of lines through one point. For a pencil through
a rational point on $E$, we compute a boundary coefficient using both
real integration and congruences at every prime. Its one-variable
continuation has magnitude comparable to $|b|^{-r/2}$. The assertion
concerns the reduced sum defined in \eqref{eq:app-trace}.

Let
\[
E:\ F(X_1,X_2,X_3)=0
\]
be a smooth plane cubic with an integral equation over $\mathbb Q$, embedded by
$|3O|$, and put
\[
r=\operatorname{rank}E(\mathbb Q).
\]
Fix a nonflex rational point $P$, represented by a primitive
integral vector $P\in\mathbb Z^3$. Set
\[
\Lambda_P=P^\perp\cap\mathbb Z^3.
\]
Choose an oriented basis $v_1,v_2$ with
$v_1\times v_2=P$, and write $B=(v_1,v_2)$.
Then
\[
\operatorname{covol}(\Lambda_P)=\|P\|_2.
\]
Primitivity of $P$ implies that
$B^T:\mathbb Z^3\to\mathbb Z^2$ is surjective.
Choose an integral linear right inverse
\[
A:\mathbb Z^2\longrightarrow\mathbb Z^3,
\qquad B^TA(m)=m.
\]

For $m\in\mathbb Z^2$, define
\[
f_m(T)=F(A(m)+TP)=a_mT^2+b_mT+c_m,
\qquad D_m=b_m^2-4a_mc_m.
\]
The cubic term vanishes because $F(P)=0$.
The coefficients $a_m,b_m,c_m$ are homogeneous
in $m$ of degrees $1,2,3$, respectively.

We exclude $m=0$, the unique projective direction
$a_m=0$, and the finitely many directions $D_m=0$.
Denote the remaining oriented frequencies by
$\mathcal M_P^\circ$.

\subsection{Finite reflection and its exact local factors}

Write $e(x)=\exp(2\pi i x)$. For $q\ge1$ let
\[
c_q(n)=\sum_{\substack{a\bmod q\\(a,q)=1}}
e(an/q),
\qquad
T_q(v)=\sum_{x\bmod q}c_q(F(x))e(v\cdot x/q).
\]
Finite Fourier orthogonality gives
\begin{equation}\label{eq:app-reflection}
\sum_{n\bmod q}T_q(Bn)e(-m\cdot n/q)
=q^2C_q(m),
\end{equation}
where
\[
C_q(m)=\sum_{\lambda\bmod q}c_q(f_m(\lambda)).
\]
Indeed, the two Fourier equations force
$B^Tx=m\bmod q$, whose solutions are exactly
$x=A(m)+\lambda P\bmod q$.

Define, initially for $\Re t>1$,
\[
\mathcal Z_m(t)=
\sum_{q\ge1}\frac{C_q(m)}{q^{1+t}}.
\]
For each prime $p$, put
\[
N_{p,k}(m)=
\#\{\lambda\bmod p^k:
f_m(\lambda)\equiv0\pmod{p^k}\},
\qquad N_{p,0}=1.
\]
The Ramanujan-sum identity gives
\[
C_{p^k}(m)=
p^k\bigl(N_{p,k}(m)-N_{p,k-1}(m)\bigr).
\]
Consequently
\begin{equation}\label{eq:app-local}
\boxed{
\mathcal Z_{m,p}(t)=
(1-p^{-t})\sum_{k\ge0}N_{p,k}(m)p^{-kt}.
}
\end{equation}
This formula includes every prime, including primes
dividing the coefficients, content, or discriminant
of $f_m$.

If $p\nmid2a_mD_m$, then
\[
\mathcal Z_{m,p}(t)
=1+\chi_{D_m}(p)p^{-t}.
\]
Thus
\begin{equation}\label{eq:app-quadratic}
\mathcal Z_m(t)
=
\frac{L(t,\chi_{D_m})}{\zeta(2t)}
\prod_{p\mid2a_mD_m}\mathcal E_{p,m}(t),
\end{equation}
where $\chi_{D_m}$ is the primitive character of the quadratic
field for nonsquare $D_m$, and the trivial character of conductor
one for square $D_m$. Writing $X=p^{-t}$, we have
\[
\mathcal E_{p,m}(t)=
\mathcal Z_{m,p}(t)\frac{1-\chi_{D_m}(p)X}{1-X^2}.
\]
In particular the baseline factor at a ramified prime is $1-X^2$.
For square $D_m$ the global baseline is $\zeta(t)/\zeta(2t)$.

\subsection{One-sided control at the finite boundary}

We take $t\downarrow0$ along the positive reals. Local correction
factors can contain $(1+p^{-t})^{-1}$, with poles at imaginary points
approaching zero as $p$ grows. We use compact subsets of $\Re t>0$
for continuation and separate estimates on the positive real axis for
the boundary limit.

Here is a uniform arithmetic bound.

For a separable quadratic $f(T)=aT^2+bT+c$
with discriminant $D\ne0$, every root modulo
$p^k$ satisfies
\[
(2aT+b)^2\equiv D\pmod{p^k}.
\]
The elementary lifting bounds for quadratic
congruences give
\begin{equation}\label{eq:app-root-bound}
N_{p,k}\ll
p^{v_p(2a)+\lceil v_p(D)/2\rceil},
\end{equation}
uniformly in $k$, with an absolute implied
constant (one may take a factor $4$ after
accounting for the dyadic case).

If $D$ is nonsquare in $\mathbb Q_p$, then
$N_{p,k}=0$ once $k$ exceeds a quantity bounded
in terms of $1+v_p(2a)+v_p(D)$.

For real $0<X<1$,
\[
(1-X)\sum_{k\ge0}N_{p,k}X^k
\le\sup_{k\ge0}N_{p,k}.
\]
At a nonsplit exceptional prime the root counts eventually vanish.
For odd $p$ it suffices that $k>v_p(D_m)$, and for $p=2$ that
$k>v_2(D_m)+2$, since any higher square congruence would make the
discriminant a local square. The factor $(1-X)$ in
\eqref{eq:app-local} cancels the zero of the baseline at $X=1$.
At an inert prime the correction is $\sum_k N_{p,k}X^k$, and at a
ramified prime it is $\sum_k N_{p,k}X^k/(1+X)$. At a split prime
the correction on $0<X<1$ is bounded by $\sup_k N_{p,k}$.
Multiplying \eqref{eq:app-root-bound} and the length bounds over
$p\mid 2a_mD_m$ costs at most a fixed power of $|2a_mD_m|$.

On compact subsets of $\Re t>0$, the same finite corrections have
polynomial bounds in $\|m\|_2$, using the absolute values of the finite
sums or the geometric series in \eqref{eq:app-local}.

Suppose $D_m>0$ is nonsquare. The corresponding
primitive quadratic character is nontrivial
and even, so
\[
L(0,\chi_{D_m})=0.
\]
This follows from the Dirichlet functional equation \cite{DLMF}.

 Consequently
\begin{equation}\label{eq:app-nonsquare}
\boxed{
D_m>0,\quad D_m\notin\mathbb Q^{\times2}
\quad\Longrightarrow\quad
\mathcal Z_m(0^+)=0.
}
\end{equation}

Now suppose $D_m>0$ is a square. Let
$\lambda_+,\lambda_-\in\mathbb Q$ be the
distinct roots of $f_m$.

For each prime $p$, sufficiently deep root
lifting gives
\[
\lim_{k\to\infty}N_{p,k}(m)
=
p^{v_p(\sqrt{D_m})}
\left(
\mathbf1_{\lambda_+\in\mathbb Z_p}
+\mathbf1_{\lambda_-\in\mathbb Z_p}
\right).
\]
This remains valid at primes dividing
$2a_mD_m$: near an integral simple
$p$-adic root, the stable number of
solutions is $p^{v_p(f_m'(\lambda))}$.

The unexceptional Euler product is
$\zeta(t)/\zeta(2t)$, whose continued
value at zero is $1$.
Comparing the finitely many exceptional
factors gives
\begin{equation}\label{eq:app-split-value}
\boxed{
\mathcal Z_m(0^+)
=
|\sqrt{D_m}|
\prod_p
\frac{
\mathbf1_{\lambda_+\in\mathbb Z_p}
+\mathbf1_{\lambda_-\in\mathbb Z_p}
}{2}.
}
\end{equation}
The product differs from $1$ at finitely
many primes only. It may vanish if neither
root is integral at some prime.

\subsection{The separated real Mellin transform}

Let $W(x)=w(\|x\|_2)$ be a nonnegative,
nonzero, smooth radial function supported
in a compact annulus. Take a nonnegative, nonzero $\omega\in C_c^\infty(0,\infty)$, with
Mellin transform $M\omega(s)=\int_0^\infty \omega(y)y^{s-1}\,dy$.
We use the Mellin specialization of the delta kernel in
\cite{HB}.

Consider the separated Mellin kernel
\begin{align}
\mathcal A(v;t,b)
={}&\frac23\zeta(t)
\int_0^\infty\rho^{b-1}
\int_{\mathbb R^3}W(x)\notag\\
&\quad\times
\left[
M\omega(t-1)
-|F(x)|^{t-1}M\omega(1-t)
\right]e(-\rho v\cdot x)\,dx\,d\rho.
\label{eq:app-kernel}
\end{align}
Its physical Mellin relation, when used,
is
\[
t=\frac{2(s+b)}3.
\]
The symbol $s$ here denotes the height Mellin parameter. The trace
below is defined with independent variables $b,t$.

For clarity, at a finite modulus $q$ the finite-frequency multiplier
is $T_q(Bn)$. Its Fourier inversion and the Dirac-comb formula give,
in distributions,
\[
 \sum_{n\in\mathbb Z^2}T_q(Bn)e(-\rho Bn\cdot x)
 =\sum_{m\in\mathbb Z^2}C_q(m)\,
       \delta^{(2)}(\rho B^Tx-m/q).
\]
Since $[A_1,A_2,P]$ is unimodular, $x=A(y)+\lambda P$ has Jacobian
one. Integration in $y$ supplies $\rho^{-2}$. Rescaling
$\rho^{-1}=qh$ and the remaining line coordinate by $h$ gives the
argument $h(A(m)+LP)$ and the factor $q^{2-b}$. In the separated
modulus sum this changes $q^{-3-u}$ to $q^{-1-t}$, where $t=u+b$.
This computation fixes both local normalizations.

Homogeneity gives
\[
F(h(A(m)+LP))=h^3f_m(L).
\]
The singular real term becomes
\begin{equation}\label{eq:app-J}
\mathcal J_m(t,b)=
\int_0^\infty h^{3t-1-b}
\int_{\mathbb R}
|f_m(L)|^{t-1}
W(h(A(m)+LP))\,dL\,dh.
\end{equation}

For radial $W$, set
\[
M_w(z)=\int_0^\infty h^{z-1}w(h)\,dh.
\]
Since $w$ has compact annular support,
$M_w$ is entire. Radial integration gives
\begin{equation}\label{eq:app-radial}
\boxed{
\mathcal J_m(t,b)=
M_w(3t-b)
\int_{\mathbb R}
\|A(m)+LP\|_2^{b-3t}
|f_m(L)|^{t-1}\,dL.
}
\end{equation}

In particular, the radial exponent at
$t=0$ is $h^{-1-b}$, with no additional
gamma-factor zero.

Define the regular real term
\[
\mathcal U_m(b)=
\int_0^\infty h^{2-b}
\int_{\mathbb R}W(h(A(m)+LP))\,dL\,dh
=M_w(3-b)\int_{\mathbb R}\|A(m)+LP\|_2^{b-3}\,dL.
\]
The \emph{reduced separated trace} is
\begin{equation}\label{eq:app-trace}
\boxed{
\begin{aligned}
\mathcal T_P(t,b)
=\frac23\zeta(t)
\sum_{m\in\mathcal M_P^\circ}
\mathcal Z_m(t)
\bigl[
&M\omega(t-1)\mathcal U_m(b)\\
&-M\omega(1-t)\mathcal J_m(t,b)
\bigr].
\end{aligned}}
\end{equation}

The sum is over \emph{oriented} lattice
frequencies: $m$ and $-m$ both occur.

One may start with $1<\Re t<2$ and $\Re b<-B_0$, for
$B_0$ sufficiently large. The modulus series and the reduced
$m$-sum converge absolutely there. On compact positive-real-part
$t$-regions we use \eqref{eq:app-quadratic} for continuation, avoiding
its poles. Taking $B_0$ larger if necessary covers a path to the
interval $0<t<1/4$.

Equation \eqref{eq:app-trace} defines the
reduced separated object; no identification
with the complete physical delta transform
is made.

\subsection{Coarea residue and dominated convergence}

Suppose $D_m>0$. The roots $\lambda_\pm$
are distinct and real, and
\[
|f_m'(\lambda_\pm)|=|\sqrt{D_m}|.
\]

The one-dimensional distribution
$|f_m(L)|^{t-1}$ has a simple pole at
$t=0$, with residue
\[
2\sum_{j=\pm}
\frac{\delta_{\lambda_j}(L)}
{|\sqrt{D_m}|}.
\]
Hence
\begin{align}
\operatorname*{Res}_{t=0}\mathcal J_m(t,b)
={}&
\frac{2}{|\sqrt{D_m}|}
\sum_{j=\pm}
\int_0^\infty h^{-1-b}
W(hx_j(m))\,dh,\label{eq:app-coarea}\\
x_j(m)&=A(m)+\lambda_jP.\notag
\end{align}

If $D_m<0$, there is no real root,
and $\mathcal J_m(t,b)$ has no
$t=0$ coarea pole.

We justify taking the one-sided residue
inside the reduced frequency sum.

Put $M=\|m\|_2$, write $L=M\ell$, and set
\[
\widehat f_m(\ell)=
M^{-3}f_m(M\ell),\quad
\alpha_m=a_m/M,\quad
\Delta'_m=D_m/M^4.
\]
The coefficients of $\widehat f_m$
are bounded independently of $m$.
Because $a_m,D_m$ are nonzero integers,
\[
|\alpha_m|\gg M^{-1},
\qquad
|\Delta'_m|\ge M^{-4}.
\]

For positive discriminant, separate
neighborhoods of the two roots.
If $\ell_j$ is a nearest root to $\ell$, the other root is
at distance at least half their separation. Therefore
\[
|\widehat f_m(\ell)|
\ge \tfrac12\sqrt{\Delta'_m}\,|\ell-\ell_j|
\]
everywhere. Also $\|A(m)/M+\ell P\|_2\asymp_P 1+|\ell|$, since
its orthogonal projection has norm bounded below and its value at
$\ell=0$ is bounded above.

For negative discriminant, completing
the square gives
\[
|\widehat f_m(\ell)|
\ge\frac{|\Delta'_m|}{4|\alpha_m|}.
\]
The change of variables contributes $M^{b-2}$ in
\eqref{eq:app-radial}. For $0<t<1/4$ and $b$ sufficiently negative,
$t\int(1+|\ell|)^{-2}\min_j|\ell-\ell_j|^{t-1}\,d\ell$ is bounded
uniformly in the roots: integrate separately within distance one of
a root and on the complement. The positive- and negative-discriminant
bounds cost at most $O(M^2)$ and $O(M^4)$ respectively.

Using \eqref{eq:app-radial} and the
annular support of $W$, one obtains,
for fixed sufficiently negative real $b$,
\begin{equation}\label{eq:app-majorant}
\sup_{0<t<t_0}
\left|t\mathcal J_m(t,b)\right|
\ll_{b,P,W} M^{b+2}.
\end{equation}
The estimate is deliberately coarse;
its purpose is summability, not a sharp
frequency exponent.

The quadratic Dirichlet factors and
the exceptional local corrections admit
one-sided bounds polynomial in $M$,
by \eqref{eq:app-root-bound}, the conductor
bound, and the functional equation.
For each fixed positive nonsquare discriminant, the finite factor tends to zero.

Thus, for a sufficiently large $B_0$,
the terms of
\[
t\mathcal T_P(t,b)
\]
have a summable majorant when
$b<-B_0$.

Dominated convergence now applies to
$t\downarrow0$ along the real axis.
No uniform complex disk is required.

The $\mathcal U_m(b)$ term has no
$t=0$ pole. For $D_m>0$ nonsquare,
the zero \eqref{eq:app-nonsquare}
cancels the real coarea pole, so its
contribution to the residue is zero.
For $D_m<0$, the archimedean term
has no coarea pole.

Only positive square discriminants
survive.

Since $\zeta(0)=-1/2$, the two minus
signs from \eqref{eq:app-trace} and
$\zeta(0)$ cancel. For each oriented
split frequency, the residue is
\begin{equation}\label{eq:app-oriented}
\boxed{
\frac23M\omega(1)
\frac{\mathcal Z_m(0^+)}
{|\sqrt{D_m}|}
\sum_{j=\pm}\mathcal W_b(x_j(m)),
}
\end{equation}
where
\[
\mathcal W_b(x)=
\int_0^\infty h^{-1-b}W(hx)\,dh
=M_w(-b)\|x\|_2^b.
\]

\subsection{Exact content summation}

Choose a primitive representative $m_0$
of each split projective direction,
identifying $m_0$ with $-m_0$.
Let $\mathcal P_P$ denote these
projective directions, with the
linear and repeated-root exceptions
removed.

For a primitive split $m_0$, write
its distinct rational roots as
\[
\lambda_\pm=\frac{u_\pm}{d_\pm},
\qquad
(u_\pm,d_\pm)=1,\quad d_\pm>0.
\]

For $n\ge1$,
\[
f_{nm_0}(T)=n^3f_{m_0}(T/n),
\]
so
\[
\lambda_\pm(nm_0)=n\lambda_\pm,
\qquad
\sqrt{D_{nm_0}}=n^2\sqrt{D_{m_0}}.
\]
Consequently \eqref{eq:app-split-value}
gives
\[
\frac{\mathcal Z_{nm_0}(0^+)}
{|\sqrt{D_{nm_0}}|}
=
\prod_p
\frac{
\mathbf1_{v_p(n)\ge v_p(d_+)}
+\mathbf1_{v_p(n)\ge v_p(d_-)}
}{2}.
\]

The real roots also scale:
\[
x_\pm(nm_0)=nx_\pm(m_0),
\qquad
\mathcal W_b(nx)=n^b\mathcal W_b(x).
\]

For $\Re b<-1$, the content sum
converges absolutely and equals
\begin{align}
&\sum_{n\ge1}n^b
\prod_p
\frac{
\mathbf1_{v_p(n)\ge v_p(d_+)}
+\mathbf1_{v_p(n)\ge v_p(d_-)}
}{2}\notag\\
&\qquad=
\boxed{
\zeta(-b)
\prod_{p\mid d_+d_-}
\frac{
p^{b v_p(d_+)}
+p^{b v_p(d_-)}
}{2}.
}
\label{eq:app-content}
\end{align}

This includes all primes and every
nonprimitive content.

Let
\[
\mathfrak d_{m_0}(b)=
\prod_{p\mid d_+d_-}
\frac{
p^{b v_p(d_+)}
+p^{b v_p(d_-)}
}{2}.
\]

Because the original trace sums over
oriented frequencies, both $m_0$
and $-m_0$ contribute. Their radial
weights agree. Thus the exact reduced
residue, written as a projective sum, is
\begin{equation}\label{eq:app-final-residue}
\boxed{
\begin{aligned}
\mathcal R_P(b)
:={}&\lim_{t\downarrow0}t\mathcal T_P(t,b)\\
={}&\frac43M\omega(1)\zeta(-b)
\sum_{m_0\in\mathcal P_P}
\mathfrak d_{m_0}(b)
\sum_{j=\pm}\mathcal W_b(x_j(m_0)).
\end{aligned}}
\end{equation}

The coefficient $4/3$ includes the
factor two for the two orientations.
If the trace is defined using one
representative modulo sign from the
outset, its coefficient is $2/3$.

Equation \eqref{eq:app-final-residue}
is proved initially for $b<-B_0$.
Its right-hand side then defines the
one-variable meromorphic continuation
in $b$.

The factor $\zeta(-b)$ has its pole
at $b=-1$. Continuing past that pole
does not mean that the original
positive content series converges
for $-1<b<0$.

\subsection{The Mordell--Weil exponent}

Here is the height calculation and the convergence needed for this
last continuation. The pencil of lines through $P$ defines
a morphism
\[
\pi_P:E\longrightarrow\mathbb P^1
\]
of degree two. Its rational split
fibres, away from finitely many branch
directions, are parametrized by
\[
Q\longmapsto\{Q,-P-Q\},
\qquad Q\in E(\mathbb Q).
\]

The associated degree-two height satisfies
\[
h(\pi_P(Q))
=
2\widehat h(Q)
+O_P(\sqrt{\widehat h(Q)}+1).
\]
Mordell--Weil lattice counting therefore
gives, for $r>0$,
\begin{equation}\label{eq:app-pencil-count}
\#\{m_0\in\mathcal P_P:
\|m_0\|_2\le X\}
\asymp_{E,P}(\log X)^{r/2}.
\end{equation}

We now bound the weights in
\eqref{eq:app-final-residue}.

Write $M=\|m_0\|_2$.
Since $a_m,b_m,c_m$ have degrees
$1,2,3$ and $a_m$ is a nonzero integer,
the quadratic formula gives
\[
|\lambda_\pm|\ll_P M^2.
\]
The rational-root theorem gives
\[
d_\pm\mid |a_{m_0}|,
\qquad d_\pm\ll_P M.
\]

Because $B^Tx_j(m_0)=m_0$,
\[
\|x_j(m_0)\|_2\gg_P M.
\]
The root bound gives
\[
\|x_j(m_0)\|_2\ll_P M^2.
\]

For real $b<0$,
\[
0<\mathfrak d_{m_0}(b)\le1.
\]
By the arithmetic--geometric mean
inequality,
\[
\mathfrak d_{m_0}(b)
\ge(d_+d_-)^{b/2}
\gg_P M^b.
\]

Since $W$ is radial and nonnegative,
$M_w(-b)>0$ for real $b$.
Uniformly for $b<0$ sufficiently close
to zero,
\[
c_{E,P}M^{3b}
\le
\mathfrak d_{m_0}(b)
\sum_{j=\pm}\mathcal W_b(x_j(m_0))
\le
C_{E,P}M^b.
\]

Partial summation applied to
\eqref{eq:app-pencil-count} gives
\[
\sum_{m_0\in\mathcal P_P}M^{3b}
\asymp_{E,P}|b|^{-r/2},
\qquad
\sum_{m_0\in\mathcal P_P}M^b
\asymp_{E,P}|b|^{-r/2}
\]
as $b\uparrow0$ through negative
real values.

The projective sum in
\eqref{eq:app-final-residue} converges
locally uniformly for $\Re b<0$.
The factor $\zeta(-b)$ supplies a
meromorphic continuation there,
with its explicit pole at $b=-1$.

As $b\uparrow0$,
\[
\zeta(-b)\longrightarrow
\zeta(0)=-\frac12.
\]
All remaining factors in
\eqref{eq:app-final-residue} are
positive for real $b<0$.

We have proved the following result.

\medskip
\noindent\textbf{Theorem (reduced coarea--Euler rank exponent).}
{\itshape
Let $E/\mathbb Q$ be a smooth integral
plane cubic embedded by $|3O|$, and
let $P\in E(\mathbb Q)$ be nonflex.
Form the reduced separated trace
\eqref{eq:app-trace}, excluding the
zero, linear and repeated-root
directions.

For sufficiently negative real $b$,
its one-sided $t=0$ residue exists
after continuation of the complete
all-prime arithmetic Euler factors,
and is given by
\eqref{eq:app-final-residue}.

The right-hand side continues
meromorphically to $\Re b<0$,
apart from its explicit zeta pole,
and, if $r>0$, satisfies
\[
\boxed{
|\mathcal R_P(b)|
\asymp_{E,P,W,\omega}
|b|^{-r/2}
\qquad(b\uparrow0,\ b<0).
}
\]
For $r=0$ the reduced split sum is
finite and remains bounded; it may
vanish if there are no separable
split directions.
}

\medskip

The sign is also determined near zero.
Since $\zeta(-b)<0$ for $-1<b<0$,
\[
\mathcal R_P(b)<0
\]
there whenever the split sum is nonempty.

If $r>0$ is odd, the asymptotic
$|\mathcal R_P(b)|\asymp|b|^{-r/2}$
rules out a meromorphic extension of
this \emph{one-variable continued
reduced trace} across $b=0$.
Indeed, a nonzero meromorphic function
has an integral Laurent order, which
cannot be comparable to a nonintegral
half-power.

This conclusion concerns only that
one-variable continuation. It does not
assert failure of joint meromorphy for
the original two-parameter transform.

\subsection{Exact scope of the residue}

The arithmetic factor in this appendix
uses every prime. Omitting a fixed
prime set can cancel the trivial zero
of a positive nonsquare quadratic
Dirichlet $L$-function and would
invalidate the split-fibre projection.

The real residue is a coarea residue,
not the full value at $t=0$.
Positive nonsquare discriminants can
contribute to the finite part after
the arithmetic zero cancels the
archimedean pole; they do not contribute
to the $t^{-1}$ coefficient.

The omitted directions are $m=0$, $a_m=0$ and $D_m=0$.
Their contributions and the remaining contour terms are needed for
the full height problem.

The order of operations is essential. We first take
$\lim_{t\downarrow0}t\mathcal T_P(t,b)$ with $b$ sufficiently negative,
and then continue that coefficient in $b$. The behavior of the
complete height transform near its central parameters is open.


\section{A positive conormal limit and a Kummer-density obstruction}
\label{app:conormal}

Let $E/\mathbb Q$ be a smooth elliptic curve embedded
in $\mathbb P^2$ by $|3O|$, with primitive integral
homogeneous cubic equation
\[
F(X_1,X_2,X_3)=0.
\]
Fix a good prime $p\ge5$, and put
\[
n_p=\#E(\mathbb F_p)=p+1-a_p(E),\qquad
c_p=\frac{(p-1)n_p}{p^2}.
\]
Normalize additive Haar measure by $\operatorname{vol}(\mathbb Z_p^3)=1$, and write $\psi_p$ for a character of $\mathbb Q_p$ with kernel $\mathbb Z_p$. The measures on compact groups used as probability laws have mass one.

\subsection{Exact finite-level squared-shell distribution}

For $k\ge1$, set
\[
\mathcal A_{p,k}=
\{a\in(\mathbb Z/p^k)^3:
a\not\equiv0\pmod p,\ F(a)\equiv0\pmod{p^k}\}.
\]
Smooth Hensel lifting gives
\[
T_{p,k}:=\#\mathcal A_{p,k}
=(p-1)n_p p^{2(k-1)}.
\]

For $\xi\in\mathbb Z_p^3$, define
\[
S_{p^k}(\xi)=
\sum_{a\in\mathcal A_{p,k}}
\psi_p(p^{-k}a\cdot\xi).
\]
Orthogonality gives
\[
\int_{\mathbb Z_p^3}|S_{p^k}(\xi)|^2\,d\xi
=T_{p,k}.
\]
Consequently
\[
d\nu_{p,k}(\xi)=
T_{p,k}^{-1}|S_{p^k}(\xi)|^2\,d\xi
\]
is a probability measure.

Fix $l\ge1$ and assume $k\ge2l$.
For $u\in(\mathbb Z/p^l)^3$, Fourier
orthogonality yields
\begin{align}
&\int_{\mathbb Z_p^3}
\psi_p(-p^{-l}u\cdot\xi)\,d\nu_{p,k}(\xi)
\notag\\
&\quad=
\frac1{T_{p,k}}
\#\{a\in\mathcal A_{p,k}:
a+p^{k-l}u\in\mathcal A_{p,k}\}.
\label{eq:conormal-correlation}
\end{align}

Since $k\ge2l$, Taylor expansion gives
\[
F(a+p^{k-l}u)
\equiv F(a)+p^{k-l}\nabla F(a)\cdot u
\pmod{p^k}.
\]
The second condition in
\eqref{eq:conormal-correlation} is therefore
equivalent to
\[
\nabla F(a)\cdot u\equiv0\pmod{p^l}.
\]

Multiplying a primitive homogeneous representative
by a unit multiplies its gradient by the square
of that unit. Thus the condition descends to
the projective point.

Every point of $E(\mathbb Z/p^l)$ has the same
number of lifts contributing to $\mathcal A_{p,k}$.
Hence
\begin{equation}\label{eq:conormal-character}
\boxed{
\widehat\nu_{p,k}(p^{-l}u)=
\frac{
\#\{P\in E(\mathbb Z/p^l):
\nabla F(P)\cdot u\equiv0\pmod{p^l}\}
}{
\#E(\mathbb Z/p^l)
}.}
\end{equation}

In particular the answer has stabilized at level $2l$.

Let $P_l$ be uniform on $E(\mathbb Z/p^l)$,
and let $\lambda_l$ be independently uniform
on $\mathbb Z/p^l$.
Character orthogonality shows that
\eqref{eq:conormal-character} is precisely
the characteristic function of
\[
\lambda_l\nabla F(P_l)\pmod{p^l}.
\]

Thus the reduction of $\nu_{p,k}$ modulo $p^l$
is, for every $k\ge2l$, exactly the law
\begin{equation}\label{eq:finite-conormal-law}
\boxed{
(\lambda_l,P_l)\longmapsto
\lambda_l\nabla F(P_l)\pmod{p^l}.}
\end{equation}

In particular, for
$v\in(\mathbb Z/p^l)^3$,
\[
\nu_{p,k}(v+p^l\mathbb Z_p^3)
=
\frac{
\#\{(P,\lambda):
\lambda\nabla F(P)\equiv v\pmod{p^l}\}
}{
p^l\#E(\mathbb Z/p^l)
}.
\]
Here $P$ and $\lambda$ range over
$E(\mathbb Z/p^l)$ and $\mathbb Z/p^l$.

\subsection{The central positive spectral limit}

Let $\mathcal U_{p,k}\subset\mathbb Q_p^3$
be the union of the cosets
\[
p^{-k}a+\mathbb Z_p^3,
\qquad a\in\mathcal A_{p,k}.
\]
These denominator shells are pairwise disjoint.

For real $w>1$, define the $L^2$ function
\[
f_{p,w}=
\mathbf1_{\mathbb Z_p^3}
+\sum_{k\ge1}p^{-kw}\mathbf1_{\mathcal U_{p,k}}.
\]
Its Fourier transform is supported on
$\mathbb Z_p^3$.

Put $r_w=p^{2-2w}$ and define the positive
finite measure
\[
d\mu_{p,w}(\xi)=
(1-r_w)|\widehat f_{p,w}(\xi)|^2\,d\xi.
\]

Disjointness of the physical shells and
Plancherel give
\[
\|\widehat f_{p,w}\|_2^2
=
1+\sum_{k\ge1}p^{-2kw}T_{p,k}
=
1+c_p\frac{r_w}{1-r_w}.
\]
Therefore
\begin{equation}\label{eq:conormal-mass}
\mu_{p,w}(\mathbb Z_p^3)
=1-r_w+c_pr_w
\longrightarrow c_p.
\end{equation}

Fix a character of conductor at most $p^l$.
Cross-correlations between distinct denominator
shells vanish once their largest denominator
exceeds $p^l$.
The remaining cross-correlations and the finitely
many low shells contribute zero after multiplication
by $1-r_w$.

For $k\ge2l$, the normalized diagonal-shell
characteristic function is exactly
\eqref{eq:conormal-character}.
The geometric tail of the diagonal-shell masses
tends to $c_p$.

It follows that every finite-level characteristic
function converges to
$c_p$ times the characteristic function of
$\lambda\nabla F(P)$, where
\[
\lambda\sim\operatorname{Haar}_{\mathbb Z_p},
\qquad
P\sim\operatorname{Haar}_{E(\mathbb Q_p)}
\]
are independent probability variables.

As $\mathbb Z_p^3$ is compact, finite-level
Fourier convergence and \eqref{eq:conormal-mass}
prove weak convergence:
\begin{equation}\label{eq:conormal-limit}
\boxed{
\mu_{p,w}\ \xrightarrow[w\downarrow1]{\rm weak}\
c_p\,(\lambda\nabla F(P))_*
(d\lambda\,dP).}
\end{equation}

The scalar mass is
\begin{equation}\label{eq:hasse-mass}
\boxed{
c_p=
\left(1-\frac1p\right)
\frac{\#E(\mathbb F_p)}p
=
(1-p^{-1})L_p(E,1)^{-1}.}
\end{equation}

Changing a primitive homogeneous representative
of $P$ multiplies $\nabla F(P)$ by a
$p$-adic unit square. This does not change
the law of $\lambda\nabla F(P)$.

We analyze the product of the resulting local probability laws below.

\subsection{Tangent support}

Let $\gamma:E\to(\mathbb P^2)^\vee$ be
the Gauss map
\[
\gamma(P)=[\nabla F(P)].
\]
Since $E$ is smooth and $p\ge5$, the
gradient is primitive at every integral
projective point.

Write $\mu_p$ for the limiting measure in \eqref{eq:conormal-limit}. Its support is
\[
\{\lambda\nabla F(P):
P\in E(\mathbb Q_p),\
\lambda\in\mathbb Z_p\}.
\]

For a primitive rational integral frequency
$v\in\mathbb Z^3$, this gives
\begin{equation}\label{eq:rational-tangent}
\boxed{
v\in\operatorname{supp}\mu_p
\quad\Longleftrightarrow\quad
F^\vee(v)=0,}
\end{equation}
where $F^\vee$ defines the dual sextic.

Indeed, the forward implication makes
$v$ a tangent direction at a $p$-adic point,
so $F^\vee(v)=0$ in $\mathbb Q_p$ and hence
in $\mathbb Q$.

Conversely, a rational tangent line to
a smooth plane cubic has a unique geometric
tangency point: two distinct tangency points
would contribute intersection multiplicity
at least four, exceeding the degree three.
Galois invariance therefore makes its
tangency point rational.
Primitive integral lifts of this point and of its tangent direction give the same projective line over $\mathbb Q_p$. The proportionality factor between their primitive vectors is a unit.

This includes flexes. At a flex the Gauss
map ramifies, affecting deeper local masses,
but not the support characterization.

For a primitive $v$ and $l\ge1$,
\[
\mu_p(v+p^l\mathbb Z_p^3)
=
c_pp^{-l}
\frac{
\#\{P\in E(\mathbb Z/p^l):
[\nabla F(P)]=[v]\}
}{
\#E(\mathbb Z/p^l)
}.
\]
If $v$ is an actual tangent over $\mathbb Q_p$ whose tangency point reduces to a nonflex, the numerator is one. Indeed the Gauss map is injective on geometric points and is an immersion at that reduction, so a lift, when it exists, is unique modulo $p^l$.
Hence
\[
\boxed{
\mu_p(v+p^l\mathbb Z_p^3)
=\frac{p-1}{p^{2l+1}}.}
\]
The factor $\#E(\mathbb F_p)$ cancels
from this localized mass. The rational tangent lines themselves grow in number on the order of $(\log B)^{r/2}$ up to height $B$ when $r>0$, since \eqref{eq:chord-height} on the diagonal is $6\hat h(P)+O_E(1)$.

\subsection{The Kummer splitting test}

Let
\[
\Gamma=E(\mathbb Q),
\]
including its torsion subgroup, and fix
a prime $\ell$.

By the Mordell--Weil theorem, $\Gamma$
is finitely generated. Thus
\[
K=\mathbb Q(E[\ell],\ell^{-1}\Gamma)
\]
is a finite Galois extension, where
$\ell^{-1}\Gamma$ denotes the division
points of a finite generating set,
together with $E[\ell]$.

Consider a prime $p\nmid\ell$ of good
reduction which splits completely in $K$.

All of $E[\ell]$ reduces to
$E(\mathbb F_p)[\ell]$, giving
\[
E(\mathbb F_p)[\ell]\simeq
(\mathbb Z/\ell)^2.
\]
Moreover, every $Q\in\Gamma$ has an
$\ell$-division point defined over
$\mathbb F_p$. Hence
\[
\operatorname{red}_p(\Gamma)
\subseteq\ell E(\mathbb F_p).
\]

Consequently
\begin{equation}\label{eq:kummer-index}
\boxed{
\#\ell E(\mathbb F_p)
=\frac{\#E(\mathbb F_p)}{\ell^2}.}
\end{equation}

The Gauss map is injective on geometric
points of a smooth plane cubic in
characteristic $p\ge5$.
Among its $n_p$ rational tangent
directions, precisely $n_p/\ell^2$
belong to
\[
\gamma(\ell E(\mathbb F_p)).
\]

Thus the conditional proportion is
exactly
\begin{equation}\label{eq:kummer-density}
\boxed{
\frac{\#\gamma(\ell E(\mathbb F_p))}
{\#\gamma(E(\mathbb F_p))}
=\ell^{-2}.}
\end{equation}
For comparison, the proportion among all points of the dual projective plane over $\mathbb F_p$ is
\[
\frac{n_p}{\ell^2(p^2+p+1)}.
\]

Every globally rational tangent direction
satisfies these Kummer conditions at
every completely split good prime.

Let $\pi_K(X)$ count such primes up to $X$.
Chebotarev gives
\[
\pi_K(X)\sim
\frac1{[K:\mathbb Q]}\frac X{\log X}.
\]
Therefore the product of the exact
conditional local proportions is
\begin{equation}\label{eq:kummer-product}
\boxed{
\prod_{\substack{p\le X\\p\ {\rm split\ in}\ K}}
\ell^{-2}
=
\exp\left(
-\frac{2\log\ell}{[K:\mathbb Q]}
\frac X{\log X}(1+o(1))
\right).}
\end{equation}

Thus the closure of the rational tangent directions has measure zero for the product of the projectivized local Gauss probabilities. This remains true after conditioning on a unit scalar at every prime.

To compare scales, for fixed integer
$d\ge0$ put
\[
M_p^{(d)}=
(1-p^{-1})^d\frac{n_p}{p}.
\]
Hasse's inequality $|a_p|\le2\sqrt p$
gives
\[
\log M_p^{(d)}=O_d(p^{-1/2}).
\]
Consequently
\begin{equation}\label{eq:mass-bound}
\boxed{
\exp(-C_d\sqrt X)
\le
\prod_{\substack{p\le X\\p\ {\rm good}}}
M_p^{(d)}
\le
\exp(C_d\sqrt X)}
\end{equation}
for a suitable constant $C_d$.

Combining \eqref{eq:kummer-product} and
\eqref{eq:mass-bound}, the product of
Kummer conditional densities times
this Hasse-bounded scalar product is
at most
\[
\exp\left(
-\frac{2\log\ell}{[K:\mathbb Q]}
\frac X{\log X}(1+o(1))
+O_d(\sqrt X)
\right),
\]
which tends to zero.

Nevertheless, the tangent direction
$\gamma(O)$ is rational and belongs
to every one of the Kummer-restricted
local sets: $O=\ell O$ at every prime.

The bound concerns the measure of congruence restrictions under
independent sampling. A bound for the number of rational directions
of a given height would require control of their correlations.

\subsection{Correlating the reductions}

For the product $G=\prod_{p\notin S}E(\mathbb F_p)$, let $H$ be
the closure of the diagonal image of $\Gamma$. There is an explicit
correlated probability measure. Choose a single coset $a$ uniformly
in $\Gamma/n\Gamma$ for all primes, and independently choose $x$ with
Haar probability on $G$. The law of $a+[n]x$ is Haar probability on
$H+nG$. Indeed $[n]x$ is Haar on $nG$, and the image of
$\Gamma/n\Gamma$ in $G/nG$ equals the image of $H$ there. This image
is finite, and its cosets have equal weights.

On any finite product $G_T$ of these groups the density with respect
to uniform probability is therefore
\[
 \frac{|G_T[n]|}{|\operatorname{im}(\Gamma\to G_T/nG_T)|}
 \,\mathbf1_{\operatorname{im}(\Gamma)+nG_T}.
\]
As $n=j!$, the closed subgroups $H+nG$ decrease to $H$.
In fact a point outside $H$ is separated from it in a finite quotient
of $G$, whose exponent divides $j!$ for all large $j$. On every
finite quotient the Haar laws eventually agree with the image of Haar
on $H$, proving weak convergence. Taking Gauss images gives the
corresponding correlated law on projective tangent directions.
Every map used at the base is a rational translate of multiplication
on $E$. The correlation is supplied by the single global coset; the
independent product in \eqref{eq:kummer-product} has no such choice.

\subsection{What remains in the scalar normalization}

The Gauss line $M=\gamma^*\mathcal O(-1)$ is $\mathcal O_E(-6O)$
since the three partial derivatives are quadrics. For marked rational
points $P_1,\ldots,P_d$, the rigidified degree-zero lines
$t_{P_i}^*M\otimes M^{-1}$ represent $6P_i$ with the convention
$t_P(Q)=Q+P$. Their nonzero fibres give torus extensions. Taking their
fibre product over a common copy of $E$ gives an extension by
$\mathbb G_m^d$. This is Bloch's construction for the corresponding
lattice \cite{Bloch}. Its normalized count at a good prime is
$c_{p,d}=(1-p^{-1})^d n_p/p$: every base point has a split torus
torsor with $(p-1)^d$ rational points. The count is the same for
dependent marks. For example, repeating a point produces a trivial
line after changing the basis of the torus characters.

There is a simple interpolation of this scalar:
\[
 \mathcal M_d^S(z)=\prod_{p\notin S}\bigl(1+(c_{p,d}-1)p^{-z}\bigr),
 \qquad \Re z>1/2.
\]
Here $S$ contains the exceptional primes of the model. The factors
are positive for real $z\ge0$. For fixed $d$, its continuation is
\[
 \mathcal M_d^S(z)=
 \frac{U_d^S(z)}
 {L^S(E,1+z)\zeta^S(1+z)^{d-1}\zeta^S(1+2z)},\qquad U_d^S(0)=1.
\]
Indeed divide each local factor by
\[
 (1-a_pp^{-1-z}+p^{-1-2z})(1-p^{-1-z})^{d-1}(1-p^{-1-2z}).
\]
The first-order trace terms agree and the two terms in $p^{-1-2z}$
cancel in the denominator. Hasse's bound makes the quotient
$1+O_{d,\delta}(p^{-3/2+3\delta})$ for large $p$ and $\Re z\ge-\delta$,
with $0<\delta<1/6$. At $z=0$ each quotient is exactly one.
Thus $\operatorname{ord}_{z=0}\mathcal M_d^S=d-a$, with
$a=\operatorname{ord}_{s=1}L(E,s)$.

Choosing a full independent set of $d=r$ marks makes the desired
nonzero finite central limit equivalent to the BSD rank equality.
The weak convergence above concerns probability measures. Multiplying
them by any common scalar at a fixed real $z>1/2$ leaves the same
scalar in the weak limit. Neither this convergence nor Bloch's
compactness theorem for the full lattice gives a bound on that scalar
as $z$ approaches zero. The finite-level Fourier law extends the local
calculations of \cite[\S3]{CLT} in the squared norm; its global
normalization remains the open step.

\begin{thebibliography}{9}
\bibitem{CLT}
A.~Chambert-Loir and Y.~Tschinkel,
\emph{Points of bounded height on equivariant compactifications of vector
groups, II}, especially Section~3 and Lemma~3.3.
\href{https://arxiv.org/abs/math/9904008}{Authors' paper}.
\bibitem{Clay}
A.~Wiles, \emph{The Birch and Swinnerton-Dyer conjecture},
\href{https://www.claymath.org/wp-content/uploads/2022/05/birchswin.pdf}{official problem description}.
\bibitem{HB}
D. R. Heath-Brown, \emph{A new form of the circle method, and its
application to quadratic forms}, J. reine angew. Math. \textbf{481}
(1996), 149--206. \href{https://ora.ox.ac.uk/objects/uuid%3Abbd3c62f-f010-44b5-8be5-87e903fb0084}{Author's paper}.
\bibitem{DLMF}
NIST Digital Library of Mathematical Functions,
\href{https://dlmf.nist.gov/25.15}{Section 25.15, Dirichlet $L$-functions}.
\bibitem{Bloch}
S.~Bloch, \emph{A note on height pairings, Tamagawa numbers,
and the Birch and Swinnerton-Dyer conjecture}, Invent. Math.
\textbf{58} (1980), 65--76.
\href{https://virtualmath1.stanford.edu/~conrad/BSDseminar/refs/BlochBSD.pdf}{Author's paper}.
\end{thebibliography}
\end{document}
